\section{Minacce alla validità}

\subsection{Milestone 1}

\paragraph{Linkage commit--ticket.}
Il linkage avviene tramite la chiave del ticket nel messaggio di commit. I commit che
correggono bug senza citare la chiave producono falsi negativi; quelli che la citano per
altri motivi producono falsi positivi. Entrambe le direzioni distorcono il tasso di bugginess.

\paragraph{Proportion e Fix Version come Affected Version.}
Proportion introduce rumore; per ridurne il ricorso, i ticket privi di AV vengono arricchiti
con le FV tranne l'ultima, ma anche questa stima è imprecisa.

\paragraph{Release senza tag Git.}
Tre versioni Jira senza tag Git sono escluse, con possibile sottostima della bugginess
nelle classi che avrebbero dovuto comparire in quelle release.

\subsection{Milestone 2}

\paragraph{Temporal leakage strutturale.}
La 10$\times$10 CV assegna i fold casualmente senza rispettare l'ordine temporale: le
metriche tendono a essere leggermente ottimistiche.

\paragraph{Nessun tuning degli iperparametri.}
I classificatori usano i default Weka (100 alberi RF, $k=1$ IBk). Un grid search
potrebbe migliorare i risultati.

\paragraph{Classificatore interno dei metodi wrapper.}
I wrapper usano NaiveBayes internamente: il subset ottimale per NB può non esserlo per RF
o IBk (v.\ risultati M2).

\subsection{Milestone 3}

\paragraph{Training bias su A.}
BClassifierA è addestrato su A e predice su A: il bias agisce identicamente
su B+ e B, cancellandosi nel confronto.

\paragraph{Indipendenza tra NSmells e le altre feature.}
L'analisi fatta nella milestone 3.1 azzera NSmells mantenendo le altre 18 feature invariate. Nella realtà
un file privo di smell ha tipicamente anche meno LOC e meno storia di modifica.

\paragraph{Composizione dei ruleset PMD.}
Il ruleset \texttt{codestyle} usato con PMD rileva violazioni di stile che crescono linearmente con la
dimensione del file, riducendo l'indipendenza di NSmells. Studi con SonarQube potrebbero riportare risultati diversi.

\paragraph{Dipendenza dal classificatore.}
I risultati dipendono dal BClassifier scelto: un classificatore diverso (es.\ IBk) avrebbe
pesi diversi per NSmells e potrebbe produrre un impatto significativamente diverso.

\subsection{Milestone 4}

\paragraph{Non-determinismo e soggettività del prompt.}
Lo stesso prompt a un LLM può produrre C\textsubscript{X} diverse in run successivi;
la scelta di quali smell enfatizzare influenza ulteriormente il risultato. I risultati
non sono deterministicamente riproducibili né facilmente confrontabili tra esecuzioni diverse.



